% GENERATED FILE. Edit canonical lessons, then run npm run generate:books.
% canonical-source-hash: fnv1a64:d517bbca9836e390
% canonical-lessons: MR-C151-pramanik, MR-C151-dhadsi, MR-C151-bhitra, MR-C151-murkh, MR-C151-prasiddh

\chapter{Honest, Brave, Timid, Foolish, Famous}
\label{ch:mr-honest-brave-timid-foolish-famous}
\hlchaptermodality{151}

\begin{chapteropening}
\textbf{By the end of this chapter:} \emph{I can say honest, brave, timid, foolish and famous.}

Complete the last lesson of chapter 151: i can say honest, brave, timid, foolish and famous.
\end{chapteropening}

\section[prāmāṇik]{\mr{प्रामाणिक} (prāmāṇik) — honest}

\label{lesson:MR-C151-pramanik}



\begin{quote}
Before the new one: say the Marathi for lazy, then the Marathi for hard-working.
\end{quote}

\subsection*{You'll want to know: \mr{प्रामाणिक}}
\textbf{\mr{प्रामाणिक}} — \emph{prāmāṇik} — \textquotedblleft{}honest\textquotedblright{}.

\begin{grammarlens}[title={What you've built}]
One more word for this chapter.
\end{grammarlens}

\subsection*{Your turn}
\begin{itemize}
\raggedright
  \item \emph{Say it:} \emph{prāmāṇik}
  \item \emph{Say it:} \emph{prāmāṇik}, once more
  \item \emph{Say it:} say \emph{prāmāṇik}
  \item \emph{Recall:} say the Marathi for lazy, then the Marathi for hard-working, then say \emph{prāmāṇik} again
\end{itemize}

\begin{tcolorbox}[breakable,colback=teal!4,colframe=teal!35!black,title={Before you move on}]
What does \mr{प्रामाणिक} mean? (\textquotedblleft{}Honest\textquotedblright{}.) Say it once more.
\end{tcolorbox}
\section[dhāḍsī]{\mr{धाडसी} (dhāḍsī) — brave}

\label{lesson:MR-C151-dhadsi}



\begin{quote}
Before the new one: say the Marathi for hard-working, then the Marathi for honest.
\end{quote}

\subsection*{You'll want to know: \mr{धाडसी}}
\textbf{\mr{धाडसी}} — \emph{dhāḍsī} — \textquotedblleft{}brave\textquotedblright{}.

\begin{grammarlens}[title={What you've built}]
One more word for this chapter.
\end{grammarlens}

\subsection*{Your turn}
\begin{itemize}
\raggedright
  \item \emph{Say it:} \emph{dhāḍsī}
  \item \emph{Say it:} \emph{dhāḍsī}, once more
  \item \emph{Say it:} say \emph{dhāḍsī}
  \item \emph{Recall:} say the Marathi for hard-working, then the Marathi for honest, then say \emph{dhāḍsī} again
\end{itemize}

\begin{tcolorbox}[breakable,colback=teal!4,colframe=teal!35!black,title={Before you move on}]
What does \mr{धाडसी} mean? (\textquotedblleft{}Brave\textquotedblright{}.) Say it once more.
\end{tcolorbox}
\section[bhitrā]{\mr{भित्रा} (bhitrā) — timid}

\label{lesson:MR-C151-bhitra}



\begin{quote}
Before the new one: say the Marathi for honest, then the Marathi for brave.
\end{quote}

\subsection*{You'll want to know: \mr{भित्रा}}
\textbf{\mr{भित्रा}} — \emph{bhitrā} — \textquotedblleft{}timid\textquotedblright{}.

\begin{grammarlens}[title={What you've built}]
One more word for this chapter.
\end{grammarlens}

\subsection*{Your turn}
\begin{itemize}
\raggedright
  \item \emph{Say it:} \emph{bhitrā}
  \item \emph{Say it:} \emph{bhitrā}, once more
  \item \emph{Say it:} say \emph{bhitrā}
  \item \emph{Recall:} say the Marathi for honest, then the Marathi for brave, then say \emph{bhitrā} again
\end{itemize}

\begin{tcolorbox}[breakable,colback=teal!4,colframe=teal!35!black,title={Before you move on}]
What does \mr{भित्रा} mean? (\textquotedblleft{}Timid\textquotedblright{}.) Say it once more.
\end{tcolorbox}
\section[mūrkh]{\mr{मूर्ख} (mūrkh) — foolish}

\label{lesson:MR-C151-murkh}



\begin{quote}
Before the new one: say the Marathi for brave, then the Marathi for timid.
\end{quote}

\subsection*{You'll want to know: \mr{मूर्ख}}
\textbf{\mr{मूर्ख}} — \emph{mūrkh} — \textquotedblleft{}foolish\textquotedblright{}.

\begin{grammarlens}[title={What you've built}]
One more word for this chapter.
\end{grammarlens}

\subsection*{Your turn}
\begin{itemize}
\raggedright
  \item \emph{Say it:} \emph{mūrkh}
  \item \emph{Say it:} \emph{mūrkh}, once more
  \item \emph{Say it:} say \emph{mūrkh}
  \item \emph{Recall:} say the Marathi for brave, then the Marathi for timid, then say \emph{mūrkh} again
\end{itemize}

\begin{tcolorbox}[breakable,colback=teal!4,colframe=teal!35!black,title={Before you move on}]
What does \mr{मूर्ख} mean? (\textquotedblleft{}Foolish\textquotedblright{}.) Say it once more.
\end{tcolorbox}
\section[prasiddh]{\mr{प्रसिद्ध} (prasiddh) — famous}

\label{lesson:MR-C151-prasiddh}



\begin{quote}
Before the new one: say the Marathi for timid, then the Marathi for foolish.
\end{quote}

\subsection*{You'll want to know: \mr{प्रसिद्ध}}
\textbf{\mr{प्रसिद्ध}} — \emph{prasiddh} — \textquotedblleft{}famous\textquotedblright{}.

\begin{grammarlens}[title={What you've built}]
One more word for this chapter.
\end{grammarlens}

\subsection*{Your turn}
\begin{itemize}
\raggedright
  \item \emph{Say it:} \emph{prasiddh}
  \item \emph{Say it:} \emph{prasiddh}, once more
  \item \emph{Say it:} say \emph{prasiddh}
  \item \emph{Recall:} say the Marathi for timid, then the Marathi for foolish, then say \emph{prasiddh} again
\end{itemize}

\begin{tcolorbox}[breakable,colback=teal!4,colframe=teal!35!black,title={Before you move on}]
What does \mr{प्रसिद्ध} mean? (\textquotedblleft{}Famous\textquotedblright{}.) Say it once more.
\end{tcolorbox}
